\documentclass[10pt,a4paper]{article}
\usepackage[utf8]{inputenc}
\usepackage[T1]{fontenc}
\usepackage[margin=0.6in]{geometry} 
\usepackage{amsmath,amssymb}
\usepackage{hyperref}
\usepackage{xcolor}

\definecolor{darkblue}{RGB}{2, 132, 199}
\hypersetup{colorlinks=true, linkcolor=darkblue, urlcolor=darkblue}

\setlength{\parindent}{0pt}
\pagestyle{empty}

\newcommand{\cvsection}[1]{%
    \vspace{12pt}
    {\large\bfseries\MakeUppercase{#1}}
    \vspace{2pt}\hrule\vspace{5pt}
}

\begin{document}

% ==========================================
% HEADER
% ==========================================
\begin{center}
    {\LARGE \bfseries Joan Alcaide-Núñez} \\[3pt]
    \small
    Email: \href{mailto:joan.alcaide@campus.lmu.de}{joan.alcaide@campus.lmu.de} $\cdot$
    Website: \href{https://joanalnu.github.io}{joanalnu.github.io} $\cdot$
    GitHub: \href{https://github.com/joanalnu}{github.com/joanalnu}
\end{center}

% ==========================================
% EDUCATION
% ==========================================
\cvsection{Education}

\textbf{Ludwig-Maximilians-Universität (LMU) Munich} \hfill Munich, Germany \\
B.Sc. in Physics \hfill Oct 2025 -- Present \\
\textit{Core concentration in foundational and mathematical physics pathways.}

\vspace{5pt}

\textbf{German School of Barcelona} \hfill Barcelona, Spain \\
Dual High School Diploma (Germany + Spain) \hfill Sept 2017 -- May 2025 \\
\textit{Academic Achievement Award with a perfect Grade Point Average of 1.0 / 1.0.}

\vspace{5pt}

\textbf{Advanced Training:} Gravity and Black Holes (\textit{Perimeter Institute GoPhysics!}, 2025) $\cdot$ Mathematics, Statistics, and Data Science (\textit{Autonomous University of Barcelona}, 2022) $\cdot$ Algorithms and C++ Programming (\textit{UPC / Harbour Space}, 2022--2023).

% ==========================================
% RESEARCH EXPERIENCE
% ==========================================
\cvsection{Research Experience}

\textbf{Theoretical Probes of Spacetime} \hfill Remote / Independent Study \\
Independent Research Project \hfill July -- Sept 2026
\begin{itemize}
    \itemsep0em \topsep2pt \parsep1pt
    \item Derived Modified Dispersion Relations (MDR) for high-energy photons traveling through non-standard cosmological backgrounds to probe potential Planck-scale Lorentz invariance violations.
    \item Applied Hamilton-Jacobi equations to analytically calculate cumulative time-of-flight delays ($\Delta t$) for gamma-ray bursts propagating under Quintessence-based FLRW metrics.
\end{itemize}

\textbf{GRB Cosmography \& Relativistic Foundations} \hfill Merate, Italy \\
Osservatorio Astronomico di Brera (OAB-INAF) \hfill June -- July 2025
\begin{itemize}
    \itemsep0em \topsep2pt \parsep1pt
    \item Constrained global cosmological parameters ($\Omega_m, \Omega_\Lambda$) via $\chi^2$ minimisation profiles using the empirical $E_{\text{peak}}-E_{\text{iso}}$ relation of Gamma-Ray Bursts.
    \item Studied relativistic jet kinematics, special relativity corrections, and detection synergies for Einstein Telescope mock populations under the guidance of Dr. G. Ghisellini, Dr. G. Ghirlanda, and Dr. O. S. Salafia.
\end{itemize}

\textbf{PI GAMMA Project} \hfill International \\
Founder \& Project Lead -- CAPIBARA Collaboration \hfill 2024 -- Present
\begin{itemize}
    \itemsep0em \topsep2pt \parsep1pt
    \item Founded an international physics collaboration for science-driven students; leading the \href{https://capibara3.github.io/gamma/}{GAMMA Team} to evaluate scientific parameters and instrumentation requirements for a transient-observing CubeSat.
\end{itemize}

\textbf{Type Ia Supernovae: Mapping Expansion History} \hfill Barcelona, Spain \\
Institute of Space Sciences (ICE-CSIC-IEEC) \hfill June -- July 2024
\begin{itemize}
    \itemsep0em \topsep2pt \parsep1pt
    \item Modeled non-linear cosmic expansion history, solving the Friedmann acceleration equations to parameterize cosmic density metrics via light-curve parameterisation using the \texttt{sncosmo} suite. \textit{Supervisor:} Dr. L. Galbany.
\end{itemize}

\textbf{Cepheid Variable Stars Research} \hfill Budapest, Hungary \\
Research Fellow -- Youth \& Science Program \hfill Aug -- Dec 2023
\begin{itemize}
    \itemsep0em \topsep2pt \parsep1pt
    \item Reconstructed classical Hubble-Lemaître expansion constraints ($H_0 = 70.26 \pm 7.1 \, \text{km} \, \text{s}^{-1} \text{Mpc}^{-1}$) using observational data from the Konkoly Observatory. Selected for a selective 3-year funded research track.
\end{itemize}

% ==========================================
% AWARDS & HONORS
% ==========================================
\cvsection{Awards \& Honors}

\textbf{Youth \& Science Fellowship} $\cdot$ \textit{Fundació Catalunya La Pedrera} \hfill 2023 -- 2026 \\
Fully funded 3-year international research fellowship supporting consecutive competitive summer research placements.

\vspace{3pt}

\textbf{Abiturpreis (Physics Distinction)} $\cdot$ \textit{German Physical Society (DPG)} \hfill 2025 \\
National prize awarded for outstanding academic excellence and foundational focus in physics during high school.

\vspace{3pt}

\textbf{International Astronomy and Astrophysics Competition (IAAC)} \hfill 2023 -- 2025 \\
Awarded the \textbf{Silver Honour} and the \textbf{National Award for Spain} (2025), ranking in the top 7\% globally.

\vspace{3pt}

\textbf{Jugend Forscht (Youth Research Competition)} \hfill 2024 -- 2025 \\
Awarded \textbf{First Prize Iberia} across regional/national tiers; Special Award in Scientific Photography (NRW, 2025).

% ==========================================
% OUTREACH & TALKS
% ==========================================
\cvsection{Talks \& Science Outreach}

\textbf{SSEA Symposium} (2026): Presented \textit{``The CRD Instrument and the COSMOS Program.''} $\cdot$ \textbf{CosmoXarxa} (2025): Invited presentation on high-energy astrophysics metrics at the Barcelona Science Museum $\cdot$ \textbf{Workshops} (2023--2024): Designed introductory physics and cosmology seminars for high school tracks at the German School Barcelona.

% ==========================================
% WORK EXPERIENCE & SKILLS
% ==========================================
\cvsection{Employment \& Technical Skills}

\textbf{LMU Munich Assistantship:} Student Assistant at Diversity Management \& Research Funding Unit \hfill 2025 -- Present \\
\textbf{Programming \& Tools:} Python (\textit{NumPy, SciPy, Pandas, Matplotlib, Astropy}), \texttt{sncosmo}, C++, \LaTeX, Markdown, Git, Linux.

\end{document}
